% 核心结论：CUA 从结构化接口逐步转向可问责闭环系统，但每一阶段都会暴露新的系统级瓶颈。
% 图原型：原创横向时间线；左侧 2009 起点 + 5 个等高时代卡 + 底部全宽 OSWorld 证据面板。
% 证据层级：hero = 五色时代序列及第 5 阶段虚线边界；支撑 = 灰底 OSWorld 成功率锚点图。
% 能不能更少？：起点、五阶段与证据面板分别承担历史起源、范式演进和数值锚定，均不可删。
% Step ①（form A，中等档）：AI Agent / CUA 综述；TikZ；目标印刷宽度 29 cm，最小字号 7 pt。
% 画布约 29 cm × 13 cm；信息流左→右；卡片以 positioning 等距串联，面板由首尾节点锚点构造。
% 模块：灰色 origin、五个低饱和时代卡、一个浅灰 evidence strip；主线为实线短箭头，前史为灰色虚线断轴。
% ASCII： [2009 小卡] - -//- - - - -┐ [时代1] -> [时代2] -> [时代3] -> [时代4] -> [时代5 虚线]
%                               └----↑（前史轴接入时代2底部）
%        └──────────────────── OSWorld 成功率锚点（全宽折线/散点 + human 虚线）────────────────────┘
% 空间预防：时代卡同高同顶；长文案分四个语义段；BJudge 标签上置，OpenCUA 标签左下置，避开 human 基线。
% Pre-flight：卡间箭头均为 node.east→node.west；无跨区 rail；证据面板四周留白均由 panel anchors 控制。

\documentclass[border=8pt]{standalone}
\usepackage{fontspec}
\usepackage{xeCJK}
\usepackage{xcolor}
\usepackage{tikz}
\setmainfont{Arial}
\setsansfont{Arial}
% PingFang is installed as a macOS downloadable TTC; XeTeX needs its path.
\setCJKmainfont{PingFang SC}[
  Path=/System/Library/AssetsV2/com_apple_MobileAsset_Font8/86ba2c91f017a3749571a82f2c6d890ac7ffb2fb.asset/AssetData/,
  UprightFont=PingFang.ttc,
  BoldFont=PingFang.ttc,
  UprightFeatures={FontIndex=3},
  BoldFeatures={FontIndex=11}
]
\setCJKsansfont{PingFang SC}[
  Path=/System/Library/AssetsV2/com_apple_MobileAsset_Font8/86ba2c91f017a3749571a82f2c6d890ac7ffb2fb.asset/AssetData/,
  UprightFont=PingFang.ttc,
  BoldFont=PingFang.ttc,
  UprightFeatures={FontIndex=3},
  BoldFeatures={FontIndex=11}
]
\usetikzlibrary{positioning,calc,arrows.meta,bending}

% Low-saturation timeline palette.
\definecolor{textMain}{HTML}{28323C}
\definecolor{textMuted}{HTML}{66717C}
\definecolor{axisGrey}{HTML}{7A858F}
\definecolor{originFill}{HTML}{F0F2F3}
\definecolor{originBorder}{HTML}{98A1A8}
\definecolor{eraOneFill}{HTML}{E7EDF4}
\definecolor{eraOneBorder}{HTML}{75879A}
\definecolor{eraTwoFill}{HTML}{E5EFEA}
\definecolor{eraTwoBorder}{HTML}{718C7D}
\definecolor{eraThreeFill}{HTML}{EEE9F3}
\definecolor{eraThreeBorder}{HTML}{887A99}
\definecolor{eraFourFill}{HTML}{F3EDE2}
\definecolor{eraFourBorder}{HTML}{9B896B}
\definecolor{eraFiveFill}{HTML}{F2E7E9}
\definecolor{eraFiveBorder}{HTML}{987B82}
\definecolor{panelFill}{HTML}{F5F6F7}
\definecolor{panelBorder}{HTML}{B6BDC3}
\definecolor{gridGrey}{HTML}{D9DDE0}
\definecolor{pointLine}{HTML}{587A94}

\newcommand{\eracontent}[4]{%
  \begin{minipage}{4.05cm}\raggedright
    {\fontsize{8.7}{10.5}\selectfont\bfseries #1}\par\vspace{4pt}
    {\fontsize{7.2}{9.6}\selectfont #2}\par\vspace{5pt}
    {\fontsize{7.2}{9.0}\selectfont\color{textMuted}#3}\par\vspace{5pt}
    {\fontsize{7.0}{8.6}\selectfont #4}
  \end{minipage}%
}

\begin{document}
\sffamily
\begin{tikzpicture}[
  every node/.style={text=textMain},
  era card/.style={
    rectangle, minimum width=4.45cm, minimum height=5.65cm,
    inner xsep=0.20cm, inner ysep=0.20cm, align=left,
    line width=0.85pt
  },
  origin card/.style={
    rectangle, minimum width=2.70cm, minimum height=2.75cm,
    text width=2.30cm, inner sep=0.18cm, align=left,
    fill=originFill, draw=originBorder, line width=0.75pt
  },
  arrow short/.style={
    -{Stealth[length=3pt 1.0,width'=0pt 0.5,sep=0pt 0.3]},
    line width=0.8pt, shorten >=1pt, shorten <=0pt,
    color=axisGrey
  },
  point label/.style={font=\fontsize{7.2}{8.8}\selectfont, text=textMain, align=left},
  tick label/.style={font=\fontsize{7.0}{8.2}\selectfont, text=textMuted}
]

% ---------------------------------------------------------------------
% Top row: origin + five eras. All placements flow from node anchors.
% ---------------------------------------------------------------------
\node[origin card, minimum width=2.70cm, minimum height=2.75cm] (origin) {%
  {\fontsize{8.4}{10.0}\selectfont\bfseries 2009 · Sikuli (UIST)}\par\vspace{5pt}
  {\fontsize{7.1}{9.5}\selectfont pixels-in,\\keyboard-mouse-out\\范式起点;\\纯模板匹配,\\无语义泛化}%
};

\node[era card, minimum width=4.45cm, minimum height=5.65cm,
      fill=eraOneFill, draw=eraOneBorder, right=1.08cm of origin] (era1) {%
  \eracontent
    {2017–2023\\结构化接口}
    {DOM·AXTree + element action\\+ self-hosted Web}
    {瓶颈:绑定网页结构,难迁移\\canvas·Mobile·Desktop}
    {WebArena}%
};

\node[era card, minimum width=4.45cm, minimum height=5.65cm,
      fill=eraTwoFill, draw=eraTwoBorder, right=0.80cm of era1] (era2) {%
  \eracontent
    {2023–2024\\Screenshot-native}
    {高分辨率视觉 + coordinate grounding}
    {瓶颈:局部 grounding 与\\长程成功脱节}
    {CogAgent · OmniParser}%
};

\node[era card, minimum width=4.45cm, minimum height=5.65cm,
      fill=eraThreeFill, draw=eraThreeBorder, right=0.80cm of era2] (era3) {%
  \eracontent
    {2024–2025\\Agent-system 化}
    {grounder·planner·memory\\·critic·tool router\\模块化}
    {瓶颈:误差级联、\\成本上升、状态所有权不清}
    {Agent S2}%
};

\node[era card, minimum width=4.45cm, minimum height=5.65cm,
      fill=eraFourFill, draw=eraFourBorder, right=0.80cm of era3] (era4) {%
  \eracontent
    {2025–2026上半年\\闭环学习}
    {task·state·verifier 共生成 + online RL + environment factory}
    {瓶颈:task validity、\\rollout 吞吐、verifier 偏差}
    {EvoCUA · DreamGym}%
};

\node[era card, minimum width=4.45cm, minimum height=5.65cm,
      fill=eraFiveFill, draw=eraFiveBorder,
      dash pattern=on 5pt off 3pt, right=0.80cm of era4] (era5) {%
  \eracontent
    {2026-07\\可问责系统(萌芽)}
    {belief provenance + explicit task state + semantic action + oversight}
    {瓶颈:跨层因果证据、\\安全边界未闭合}
    {GUIStateBelief · Tactile}%
};

% Pre-LLM span: dashed gray axis routes under Era 1 and enters Era 2.
\coordinate (preaxisLevel) at ([yshift=-0.30cm]era1.south);
\coordinate (preaxisExit) at ([xshift=0.18cm]origin.east);
\coordinate (preaxisDown) at (preaxisExit |- preaxisLevel);
\coordinate (preaxisUnderEraTwo) at (era2.south |- preaxisLevel);
\draw[axisGrey, line width=0.8pt, densely dashed]
  (origin.east) -- (preaxisExit) -- (preaxisDown) -- (preaxisUnderEraTwo);
\draw[arrow short, densely dashed] (preaxisUnderEraTwo) -- (era2.south);
\coordinate (prebreak) at ($(preaxisDown)!0.50!(preaxisUnderEraTwo)$);
\draw[white, line width=2.6pt] ([xshift=-2pt]prebreak) -- ([xshift=2pt]prebreak);
\draw[axisGrey, line width=0.65pt]
  ([xshift=-3pt,yshift=-3pt]prebreak) -- ([xshift=1pt,yshift=3pt]prebreak)
  ([xshift=1pt,yshift=-3pt]prebreak) -- ([xshift=5pt,yshift=3pt]prebreak);

% Adjacent-era progression arrows; every gap is shorter than 1.5 cm.
\draw[arrow short] (era1.east) -- (era2.west);
\draw[arrow short] (era2.east) -- (era3.west);
\draw[arrow short] (era3.east) -- (era4.west);
\draw[arrow short] (era4.east) -- (era5.west);

% ---------------------------------------------------------------------
% Bottom evidence strip. Its width is constructed from the outer nodes.
% ---------------------------------------------------------------------
\coordinate (panelNW) at ([yshift=-0.50cm]era1.south west -| origin.west);
\coordinate (panelNE) at ([yshift=-0.50cm]era1.south east -| era5.east);
\coordinate (panelSW) at ([yshift=-5.75cm]panelNW);
\coordinate (panelSE) at ([yshift=-5.75cm]panelNE);
\fill[panelFill] (panelSW) rectangle (panelNE);
\draw[panelBorder, line width=0.75pt] (panelSW) rectangle (panelNE);

\coordinate (panelTopMid) at ($(panelNW)!0.5!(panelNE)$);
\node[anchor=north, font=\fontsize{9.0}{11.0}\selectfont\bfseries, text=textMain]
  at ([yshift=-0.28cm]panelTopMid)
  {OSWorld 成功率锚点 — 注意:三点来自不同系统与设置,非同一受控实验,不构成单一能力曲线};

% Chart frame: x encodes month distance; y encodes Success Rate on 0--80%.
\coordinate (chartBL) at ([xshift=1.20cm,yshift=0.82cm]panelSW);
\coordinate (chartBR) at ([xshift=-0.55cm,yshift=0.82cm]panelSE);
\coordinate (chartTL) at ([yshift=3.95cm]chartBL);
\coordinate (chartTR) at ([yshift=3.95cm]chartBR);

% Horizontal grid and y ticks.
\foreach \frac/\lab in {0/0\%,0.25/20\%,0.5/40\%,0.75/60\%,1/80\%}{
  \coordinate (gridL) at ($(chartBL)!\frac!(chartTL)$);
  \coordinate (gridR) at ($(chartBR)!\frac!(chartTR)$);
  \draw[gridGrey, line width=0.35pt] (gridL) -- (gridR);
  \node[tick label, anchor=east] at ([xshift=-0.10cm]gridL) {\lab};
}
\draw[axisGrey, line width=0.55pt] (chartBL) -- (chartBR);
\draw[axisGrey, line width=0.55pt] (chartBL) -- (chartTL);
\node[tick label, anchor=south west, text=textMain]
  at ([yshift=0.12cm]chartTL) {Success Rate};
\node[tick label, anchor=north]
  at ([yshift=-0.36cm]$(chartBL)!0.5!(chartBR)$) {time};

% Month-proportional positions with equal one-month padding on both sides:
% 2024-10 = 1/14, 2025-08 = 11/14, 2025-10 = 13/14.
\coordinate (dateOneB) at ($(chartBL)!0.0714286!(chartBR)$);
\coordinate (dateOneT) at ($(chartTL)!0.0714286!(chartTR)$);
\coordinate (dateTwoB) at ($(chartBL)!0.7857143!(chartBR)$);
\coordinate (dateTwoT) at ($(chartTL)!0.7857143!(chartTR)$);
\coordinate (dateThreeB) at ($(chartBL)!0.9285714!(chartBR)$);
\coordinate (dateThreeT) at ($(chartTL)!0.9285714!(chartTR)$);

\coordinate (pointOne) at ($(dateOneB)!0.182875!(dateOneT)$);
\coordinate (pointTwo) at ($(dateTwoB)!0.5625!(dateTwoT)$);
\coordinate (pointThree) at ($(dateThreeB)!0.9075!(dateThreeT)$);
\coordinate (humanL) at ($(chartBL)!0.9045!(chartTL)$);
\coordinate (humanR) at ($(chartBR)!0.9045!(chartTR)$);

% X ticks and exact dates.
\foreach \base/\date in {dateOneB/2024-10,dateTwoB/2025-08,dateThreeB/2025-10}{
  \draw[axisGrey, line width=0.45pt] (\base) -- ([yshift=-0.09cm]\base);
  \node[tick label, anchor=north] at ([yshift=-0.12cm]\base) {\date};
}

% Human reference and observed system anchors.
\draw[axisGrey, line width=0.75pt, densely dashed] (humanL) -- (humanR);
\node[point label, anchor=south west, text=textMuted]
  at ([xshift=0.18cm,yshift=0.10cm]humanL) {human 72.36\%};

\draw[pointLine, line width=1.05pt] (pointOne) -- (pointTwo) -- (pointThree);
\foreach \p in {pointOne,pointTwo,pointThree}{
  \fill[white] (\p) circle[radius=0.15cm];
  \fill[pointLine] (\p) circle[radius=0.105cm];
}

\node[point label, anchor=south west]
  at ([xshift=0.18cm,yshift=0.48cm]pointOne)
  {GPT-4o + OS-Atlas-Base-7B, 14.63\%};
\node[point label, anchor=south east]
  at ([xshift=-0.18cm,yshift=0.22cm]pointTwo)
  {OpenCUA-72B (OSWorld-Verified, 100 steps), 45.0\%};
\node[point label, anchor=south east]
  at ([xshift=-0.12cm,yshift=0.25cm]pointThree)
  {BJudge (100 steps), 72.6\%};

\end{tikzpicture}
\end{document}
